\documentclass[10pt,letterpaper]{article}
\usepackage[utf8]{inputenc}
\usepackage[spanish]{babel}
\usepackage{enumitem}
\usepackage{hyperref}
\usepackage{xcolor}
\usepackage{fontawesome5}
\usepackage{parskip}
\usepackage[margin=0.75in]{geometry}
\usepackage{setspace}

\hypersetup{
    colorlinks=true,
    linkcolor=black,
    urlcolor=black,
    citecolor=black,
    filecolor=black
}

\setlist[itemize]{leftmargin=*,nosep,after=\vspace{0.1cm}}
\setlist[enumerate]{leftmargin=*,nosep}
\setlength{\parskip}{0.3cm}

\begin{document}

% Header
\begin{center}
    {\Huge \textbf{Bruno Camilo Henríquez Cano}} \\[0.2cm]
    {\Large Coordinador de Taller $\cdot$ Planificador de Mantenimiento $\cdot$ Gestión Operativa}
\end{center}

\vspace{0.3cm}

% Contact
\noindent
\begin{minipage}[t]{0.5\textwidth}
    \faMapMarker\ Santiago, Chile \\
    \faPhone\ +56 9 2815 6161 \\
    \faEnvelope\ \href{mailto:bruno.henriquez.1993@gmail.com}{bruno.henriquez.1993@gmail.com}
\end{minipage}
\begin{minipage}[t]{0.5\textwidth}
    \faLinkedin\ \href{https://linkedin.com/in/bruno-henriquezcano}{linkedin.com/in/bruno-henriquezcano} \\
    
    \faGlobe\ \href{https://bchc.tech}{bchc.tech} \\
    \faCar\ Licencia de Conducir Clase B
\end{minipage}

\vspace{0.4cm}

% Professional Profile
\noindent{\Large \textbf{PERFIL PROFESIONAL}}

\vspace{0.2cm}

\noindent Profesional con más de 7 años de experiencia en coordinación de talleres de mantenimiento y planificación técnica para flotas de equipos industriales. Especializado en gestión de operaciones de taller, control de inventario de repuestos, administración presupuestaria y coordinación de equipos técnicos. Experiencia demostrada gestionando presupuestos superiores a \$2M CAD (1.500M CLP) mensuales, implementando sistemas de control de costos y optimizando procesos operativos. Capacidad técnica sólida para entender problemáticas de terreno y traducirlas en soluciones prácticas y planificación efectiva.

\vspace{0.4cm}

% Work Experience
\noindent{\Large \textbf{EXPERIENCIA LABORAL}}

\vspace{0.3cm}

\noindent \textbf{Secretario de Planificación Técnica} \\
\textit{Consorcio Acciona Ossa Pizzarotti (Codelco) — Calama — Diciembre 2023 – Enero 2025}
\begin{itemize}
    \item \textbf{Gestión de Taller:} Coordiné las operaciones diarias del taller de mantenimiento para una flota de +100 equipos mineros pesados, asegurando disponibilidad operativa y cumplimiento de programas de mantenimiento.
    \item \textbf{Control Presupuestario:} Gestioné y reporté mensualmente gastos de suministros del taller para un presupuesto promedio de \$2M CAD (1.500M CLP), implementando controles que redujeron desviaciones presupuestarias.
    \item \textbf{Planificación de Mantenimiento:} Desarrollé y ejecuté planes de mantenimiento preventivo y correctivo, priorizando intervenciones según criticidad operativa y disponibilidad de recursos.
    \item \textbf{Gestión de Repuestos:} Coordiné con proveedores la adquisición y disponibilidad de repuestos críticos, estableciendo niveles de stock mínimo y máximo para componentes prioritarios.
    \item \textbf{Reportes y KPIs:} Generé reportes diarios de disponibilidad de equipos (OEE), tiempos de reparación (MTTR) y confiabilidad (MTBF), proporcionando visibilidad operativa a la supervisión.
    \item \textbf{Coordinación de Equipos:} Lideré la comunicación entre taller, operaciones y bodega de repuestos, asegurando flujo eficiente de información y resolución ágil de problemas.
    \item \textbf{Uso de SAP-PM:} Administré órdenes de trabajo, registros de mantenimiento y trazabilidad de componentes en sistema SAP-PM.
\end{itemize}

\vspace{0.2cm}

\noindent \textbf{Coordinador de Sistemas y Análisis de Flota} \\
\textit{STP Santiago S.A — Maipú — Junio 2021 – Agosto 2023}
\begin{itemize}
    \item \textbf{Coordinación Operativa:} Coordiné el mantenimiento de una flota de +400 buses de transporte público, gestionando prioridades de taller y disponibilidad de vehículos.
    \item \textbf{Gestión de Información:} Centralicé datos de mantenimiento dispersos en múltiples sistemas, creando una base de datos unificada que mejoró la trazabilidad y el análisis de tendencias.
    \item \textbf{Optimización de Procesos:} Automaticé la generación de reportes operativos, reduciendo el tiempo de procesamiento en 30\% y liberando recursos para tareas de mayor valor.
    \item \textbf{Análisis de Desempeño:} Desarrollé dashboards visuales en Power BI para seguimiento de KPIs de mantenimiento, facilitando la toma de decisiones de la gerencia.
    \item \textbf{Mejora Continua:} Identifiqué oportunidades de mejora en procesos de taller y propuse soluciones que resultaron en un aumento del 15\% en disponibilidad de flota.
\end{itemize}

\vspace{0.2cm}

\noindent \textbf{Técnico Automotriz Independiente} \\
\textit{Autónomo — Paine, Chile — 2014 – 2020}
\begin{itemize}
    \item \textbf{Gestión Integral de Taller:} Administré de manera autónoma un taller mecánico, gestionando clientes, proveedores, inventario de repuestos y control de costos.
    \item \textbf{Atención de Clientes:} Atendí más de 30 clientes mensuales, diagnosticando fallas, cotizando reparaciones y asegurando satisfacción mediante comunicación transparente.
    \item \textbf{Gestión de Repuestos:} Negocié con proveedores condiciones comerciales favorables y mantuve inventario de repuestos de rotación rápida.
    \item \textbf{Control de Calidad:} Implementé procedimientos de verificación post-reparación que aseguraron un 95\% de satisfacción del cliente.
\end{itemize}

\vspace{0.4cm}

% Management Skills
\noindent{\Large \textbf{HABILIDADES DE GESTIÓN Y COORDINACIÓN}}
\vspace{0.3cm}

\noindent
\begin{minipage}[t]{0.48\textwidth}
\textbf{Gestión de Taller:}
\begin{itemize}
    \item Planificación de Mantenimiento (Preventivo/Correctivo)
    \item Coordinación de Equipos Técnicos
    \item Priorización de Órdenes de Trabajo
    \item Gestión de Proveedores
    \item Control de Calidad de Reparaciones
\end{itemize}

\vspace{0.2cm}

\textbf{Control Administrativo:}
\begin{itemize}
    \item Gestión Presupuestaria (\$2M CAD/mes)
    \item Control de Costos de Mantenimiento
    \item Gestión de Inventario de Repuestos
    \item Reportes Ejecutivos y KPIs
    \item Administración de Recursos
\end{itemize}
\end{minipage}
\hfill
\begin{minipage}[t]{0.48\textwidth}
\textbf{Herramientas de Gestión:}
\begin{itemize}
    \item SAP-PM (Planificación de Mantenimiento)
    \item Microsoft Excel (análisis y reportes)
    \item Power BI (dashboards operativos)
    \item Sistemas de Control de Inventario
    \item Software de Gestión de Órdenes de Trabajo
\end{itemize}

\vspace{0.2cm}

\textbf{KPIs de Mantenimiento:}
\begin{itemize}
    \item OEE (Eficiencia Global de Equipos)
    \item MTBF (Tiempo Medio Entre Fallas)
    \item MTTR (Tiempo Medio de Reparación)
    \item Disponibilidad de Flota
    \item Cumplimiento de Plan de Mantenimiento
\end{itemize}
\end{minipage}

\vspace{0.3cm}

% Technical Knowledge
\vspace{0.4cm}
\noindent{\Large \textbf{CONOCIMIENTOS TÉCNICOS}}

\vspace{0.3cm}

\begin{itemize}
    \item \textbf{Diagnóstico Técnico:} Capacidad de entender y diagnosticar fallas en sistemas mecánicos, hidráulicos y eléctricos, permitiendo comunicación efectiva con el equipo técnico.
    \item \textbf{Planificación Técnica:} Conocimiento de procedimientos de mantenimiento, frecuencias de intervención y criticidad de componentes.
    \item \textbf{Gestión de Repuestos:} Identificación de partes críticas, tiempos de aprovisionamiento y estrategias de stock.
    \item \textbf{Automatización:} Desarrollo de scripts y herramientas para automatizar reportes y mejorar eficiencia administrativa (Python, VBA).
\end{itemize}

\vspace{0.4cm}

% Education
\noindent{\Large \textbf{EDUCACIÓN}}

\vspace{0.3cm}

\noindent \textbf{Ingeniería en Maquinaria, Sistemas Automotrices y Electrónicos} \\
\textit{INACAP — Marzo 2014 – Octubre 2019}

\vspace{0.2cm}

\noindent \textbf{Ingeniería Civil en Computación (2 años completados)} \\
\textit{Universidad Tecnológica Metropolitana — Marzo 2012 – Marzo 2014}

\vspace{0.4cm}

% Languages
\noindent{\Large \textbf{IDIOMAS}}

\vspace{0.3cm}

\noindent \textbf{Español:} Nativo \\
\noindent \textbf{Inglés:} Intermedio (lectura técnica y comunicación profesional)

\end{document}
